% solutions/thm-a3-block-gibbs.tex
% Standalone proof dossier; it also lifts into main.tex as a subfile.
\documentclass[../main.tex]{subfiles}
\begin{document}

% === SOLUTION HEADER (proof provenance) =====================================
%   ledger-node : thm:a3-block-gibbs
%   refines     : thm:a3-block-gibbs
%   bounded_by  : obs:marginals-not-joint
%   evidence_run: none
%   checked_by  : agent
%   author      : /root/a3
%   reviewer    : /root/cross_review
%   review      : research/reviews/2026-08-21-a3-independent-agent-audit.md
%   review level: independent agent; no human or Lean review is claimed
%   date        : 2026-08-21
% ============================================================================

\section*{Solution: block-Gibbs tensorization and maximal correlation}
\label{sec:sol-thm-a3-block-gibbs}

\begin{theorem}
\label{thm:sol-a3-block-gibbs}
Let $X=(X_1,\ldots,X_m)\sim\pi$, with at least one nondegenerate block, and define
\[
 \gamma_{\mathrm{blk}}(\pi)
 =\inf_{\Var_\pi(f)>0}
   \frac{\sum_{i=1}^m\E_\pi[\Var(f\mid X_{-i})]}{\Var_\pi(f)}.
\]
Suppose that, for almost every $x_{-i}$, the conditional law of $X_i$ satisfies the weighted
Poincar\'e inequality
\[
 \Var_{\pi_i(\cdot\mid x_{-i})}(g)
 \le C_i\int
 \left\langle A_i(x)\nabla_i g,\nabla_i g\right\rangle
 \dd\pi_i(\cdot\mid x_{-i}),
\]
where each $C_i<\infty$ is deterministic and the displayed conditional energies are measurable.
Then, whenever $\gamma_{\mathrm{blk}}(\pi)>0$,
\[
 \Var_\pi(f)
 \le\frac{\max_iC_i}{\gamma_{\mathrm{blk}}(\pi)}
 \int\sum_{i=1}^m
 \left\langle A_i(x)\nabla_i f,\nabla_i f\right\rangle\dd\pi(x)
\]
on the joint form domain.  A product law has $\gamma_{\mathrm{blk}}=1$.  For exactly two
blocks $X,Y$,
\[
 \gamma_{\mathrm{blk}}=1-\rho_{\max}(X,Y),
\]
where $\rho_{\max}$ is the Hirschfeld--Gebelein--R\'enyi maximal correlation.  This equality
also holds when the norm is not attained; in particular, a nonconstant common factor gives
$\rho_{\max}=1$ and $\gamma_{\mathrm{blk}}=0$.
\end{theorem}

\begin{proof}
The definition of $\gamma_{\mathrm{blk}}$ gives
\[
 \gamma_{\mathrm{blk}}\Var_\pi(f)
 \le\sum_i\E_\pi\Var(f\mid X_{-i}).
\]
Applying the conditional inequality block by block, integrating in $x_{-i}$, and bounding every
$C_i$ by $\max_iC_i$ proves the joint inequality.

If the blocks are independent, the Efron--Stein inequality gives
\[
 \Var_\pi(f)\le\sum_i\E_\pi\Var(f\mid X_{-i}),
\]
so $\gamma_{\mathrm{blk}}\ge1$.  A nonconstant function of one nondegenerate block has quotient
exactly one, proving the reverse inequality.

Now take two blocks and work in the real Hilbert space $\mathcal H=L^2_0(\pi)$.  Let $P_X$ and
$P_Y$ be the orthogonal conditional-expectation projections onto the centered $X$- and
$Y$-measurable subspaces.  Conditional variance gives, for $f\in\mathcal H$,
\begin{equation}
\label{eq:sol-a3-projection-form}
 \E\Var(f\mid X)+\E\Var(f\mid Y)
 =\langle f,(2I-P_X-P_Y)f\rangle.
\end{equation}
The HGR coefficient is
\[
 \rho_{\max}(X,Y)
 =\sup\left\{|\E[a(X)b(Y)]|:
   \E a=\E b=0,\ \E a^2=\E b^2=1\right\}.
\]
Equivalently,
$\rho_{\max}=\|P_XP_Y\|$, since $P_X$ restricted to the centered $Y$-subspace is the
conditional-expectation operator whose norm is the displayed supremum.

For completeness, the two-projections identity says that orthogonal projections $P,Q$, not both
zero, satisfy
\begin{equation}
\label{eq:sol-a3-two-projections}
 \|P+Q\|=1+\|PQ\|.
\end{equation}
One way to see it is to apply the spectral theorem to $PQP$ on $\operatorname{Ran}P$.
On the two-dimensional fiber associated with a singular value $s$ of $PQ$, the two projections
have sum with eigenvalues $1-s$ and $1+s$; the intersection and one-projection fibers give the
limiting values $2$, $1$, or $0$.  Taking the supremum, including approximate spectral fibers
when the norm is not attained, yields \eqref{eq:sol-a3-two-projections}.

The operator $P_X+P_Y$ is positive.  Hence
\[
 \inf\operatorname{spec}(2I-P_X-P_Y)
 =2-\|P_X+P_Y\|
 =1-\rho_{\max}(X,Y).
\]
Together with \eqref{eq:sol-a3-projection-form}, this is precisely the variational definition of
$\gamma_{\mathrm{blk}}$ and proves the two-block formula.  The operator-norm argument uses
approximating vectors and therefore does not assume attainment.
\end{proof}

\paragraph{Obstruction respected.}
The bound contains the explicit dependence factor $\gamma_{\mathrm{blk}}^{-1}$.  Fixed marginal
Hardy constants alone give no joint estimate: a bottleneck or common factor is recorded by
$\gamma_{\mathrm{blk}}\downarrow0$.

\end{document}
